\subsection{Baselines for Optimizability Prediction}
\label{sec:method:architecture:baselines}
RQ1 is a comparative claim and needs a reference. Two groups supply one: trivial predictors,
which fix what a given $R^2$ is worth on this label, and three published circuit models
ported to this task. A third group of standard graph-level encoders (GCN, GraphSAGE, GIN)
was considered and may be added if time permits.

\subsubsection{Trivial Predictors}
\label{sec:method:baselines:trivial}
Three predictors, each fitted on the validation split and scored on test, so that every
baseline is scored on data it has not seen. Two are constant,
\begin{equation}
    \hat{y}^{\text{mean}}_i = \bar{y}_{\text{val}},
    \qquad
    \hat{y}^{\text{med}}_i = \operatorname{med}\!\left( y_{\text{val}} \right),
    \label{eq:method:baseline-constant}
\end{equation}
and the third is an ordinary least squares fit on graph size alone,
\begin{equation}
    \hat{y}^{\text{size}}_i = \operatorname{clip}_{[0,1]}\!\left(
        \beta_0
        + \beta_1 \log_{10} \left| \mathcal{V}(G_i) \right|
        + \beta_2 \log_{10} \left| \mathcal{E}(G_i) \right| \right).
    \label{eq:method:baseline-size}
\end{equation}
Equation~\eqref{eq:method:baseline-constant} fixes the zero point
of~\eqref{eq:method:r2}. The test-set mean is the least-squares optimal constant, so any
other constant scores $R^2 \leq 0$, with equality only at
$\bar{y}_{\text{val}} = \bar{y}_{\text{test}}$. Equation~\eqref{eq:method:baseline-size} is
the weakest non-trivial hypothesis about the task, that optimizability is a function of size
alone, and any claim that the encoder reads circuit \emph{structure} has to clear it.

\subsubsection{Adapting Published Circuit Models}
\label{sec:method:baselines:published}
SynthNet \citep{chowdhury2021openabc}, HOGA \citep{deng2024hoga} and DeepGate4
\citep{deepgate4_2025} are introduced in~\ref{sec:relwork:ml4eda:baselines}. None was
published for this task, so each had to be adapted, and what follows states those
adaptations rather than the models.

\paragraph{Shared porting procedure.}
Each paper supplies an encoder $g_\phi$ carrying an AIG to node states
$\{ h_v \}_{v \in \mathcal{V}(G)}$, together with a head supervising that paper's own
target. The encoder is reused unmodified and the head is replaced, so every port has the
form
\begin{equation}
    \hat{y}(G) = \operatorname{Sigmoid}\!\left(
        \operatorname{MLP}\!\left(
            \pi\!\left( \{ h_v \}_{v \in \mathcal{V}(G)} \right)
        \right) \right),
    \label{eq:method:baseline-port}
\end{equation}
with $\pi$ a permutation-invariant readout and the terminal sigmoid matching the $[0,1]$
range of the target of~\ref{sec:prelim:formalization}. Only $\pi$, the head and the input
adapter differ across the three, and the trunk of~\eqref{eq:method:baseline-port} is in
every case the published one.

Whether $\pi$ is inherited or chosen differs by model, and the distinction bears on what
each comparison licenses. SynthNet already predicts a graph-level quantity, so $\pi$ is its
own concatenated maximum and mean pooling. HOGA classifies nodes and DeepGate4 pools over
cones; neither defines a whole-circuit readout, so for both $\pi$ is mean pooling and is
this work's choice, not the paper's.

\paragraph{Adaptations required by a graph-level target.}
Three, each a consequence of the target being one scalar per circuit under one fixed script:

\begin{itemize}
    \item \textbf{The recipe branch is removed.} SynthNet predicts $\hat{y}(G, s)$,
    conditioning on a synthesis recipe $s$ through a separate encoder. Training here fixes
    the script, so $s$ is constant across the corpus and that branch emits one vector for
    every sample, contributing a bias to the head and nothing else. It is dropped rather
    than left to learn that constant.
    \item \textbf{Batches are bounded by nodes, not graphs.} A graph-level readout requires
    all of $\mathcal{V}(G)$ in one forward pass, so a graph cannot be split across batches,
    and neither published batch size transfers: HOGA minibatches nodes drawn from one
    graph, DeepGate4 minibatches cones. Batches are assembled to a node budget $B$ and
    paired with an accumulation count $a$, and it is the product that matters. One graph
    contributes one label regardless of its size, so the number of losses averaged per
    update is $a \cdot \mathbb{E}[\,B / |\mathcal{V}(G)|\,]$, held near the primary
    model's value.
    Budget and accumulation are tuned and reported as a pair.
    \item \textbf{Activations are recomputed rather than retained.} DeepGate4 attends over a
    virtual edge set $\bar{\mathcal{E}}_k = \{(u,v) : \operatorname{dist}(u,v) \leq k\}$,
    which at the published radius $k = 8$ measures roughly $180$ edges per circuit node
    against an AIG fanin of at most two. Each of its attention layers materialises a message
    tensor over $\bar{\mathcal{E}}_k$, so recomputing them during the backward pass, which
    leaves outputs and gradients unchanged, is what makes the baseline runnable at all.
\end{itemize}

\paragraph{Corrections to the released implementations.}
Two departures are corrections rather than adaptations, and are stated because they change
published behaviour. HOGA's released attention normalises its scores over the head axis
rather than the key axis, so the weights never form a distribution over hop neighbours; the
port normalises over keys, following Eq.~5 of the paper. Its optional reverse-propagation
branch applies the forward adjacency to the reverse features, which duplicates the forward
branch rather than reversing it, and the port applies the reverse adjacency. A third
difference is a convention, not an error: OpenABC-D orients edges toward the primary
inputs and this corpus orients them the
other way, so each node summarises its fanout cone under one convention and its fanin cone
under the other. Both orientations are run and the pair reported.

\paragraph{Comparability.}
Only the data pipeline is shared, since a different split would make the comparison
meaningless. Hyperparameters, optimizer, loss and schedule take each paper's published
values. Settings a paper leaves unpublished are recorded as assumptions
in~\ref{sec:apx:baselines} rather than silently absorbed.

Scores printed in those papers are a different matter and cannot be placed beside the ones
here. OpenABC-D standardises its targets within each design $d$, so its reference spread is
the within-design sum of squares where the denominator of~\eqref{eq:method:r2} is the total,
\begin{equation}
    \underbrace{\sum_{i} \left( y_i - \bar{y} \right)^{2}}_{\text{here}}
    = \underbrace{\sum_{i} \left( y_i - \bar{y}_{d(i)} \right)^{2}}_{\text{upstream}}
    + \sum_{d} n_d \left( \bar{y}_{d} - \bar{y} \right)^{2}.
    \label{eq:method:r2-decomposition}
\end{equation}
The two differ by the between-design term, which on a design-disjoint split
(\ref{sec:method:experiment:splitting}) is the variation the task is about. Only the sign
of the score transfers (\ref{sec:results:rq1:baselinemodels}).
